\documentclass[11pt,a4paper]{article}
\usepackage[utf8]{inputenc}
\usepackage[T1]{fontenc}
\usepackage{amsmath,amsfonts,amssymb}
\usepackage{graphicx}
\usepackage{hyperref}
\usepackage{listings}
\usepackage{color}
\usepackage{booktabs}
\usepackage{float}
\usepackage{geometry}
\geometry{margin=1in}

\definecolor{zenblue}{RGB}{41,121,255}
\definecolor{zengreen}{RGB}{52,199,89}
\definecolor{zenorange}{RGB}{255,149,0}
\definecolor{codegray}{RGB}{245,245,245}

\hypersetup{
    colorlinks=true,
    linkcolor=zenblue,
    urlcolor=zenblue,
    citecolor=zenblue
}

\lstset{
    backgroundcolor=\color{codegray},
    basicstyle=\ttfamily\small,
    breaklines=true,
    captionpos=b,
    frame=single,
    numbers=left,
    numberstyle=\tiny\color{gray}
}

\title{
    \vspace{-2cm}
    \Large \textbf{Zen AI Model Family} \\
    \vspace{0.5cm}
    \Huge \textbf{Zen4-Pro} \\
    \vspace{0.3cm}
    \large Fourth-Generation Professional Model \\
    \vspace{0.5cm}
    \normalsize Technical Whitepaper v2026.01
}

\author{
    Zach Kelling\thanks{zach@lux.network} \\
    \texttt{research@hanzo.ai} \\
    \\
    Zoo Labs Foundation \\
    \texttt{foundation@zoolabs.org}
}

\date{January 2026}

\begin{document}

\maketitle

\begin{abstract}
We present \textbf{Zen4-Pro}, the flagship professional model of the Zen4 generation at 72B dense
parameters. Zen4-Pro advances the state of agentic AI through deep reasoning chains, native
agentic workflows with structured plan-execute-verify loops, and improved tool orchestration
across multi-step tasks. The model achieves 91.8\% on MMLU, 87.6\% on MATH, 93.2\% on
HumanEval, 74.3\% on GPQA Diamond, and 48.7\% on SWE-bench Verified, establishing new
state-of-the-art results for the 72B dense parameter class. Zen4-Pro is designed for
enterprise deployment scenarios requiring reliable, multi-step autonomous task execution
while maintaining practical inference efficiency on standard data-center hardware.
\end{abstract}

\tableofcontents
\newpage

\section{Introduction}

Professional-tier AI deployment presents a distinct challenge: tasks are rarely single-turn
question-answer exchanges. Software engineers debug across dozens of files; data analysts chain
transformations through complex pipelines; legal professionals trace arguments across hundreds
of documents. These workflows require a model that can sustain coherent intent, manage tool
state, and self-verify intermediate results across many sequential steps.

Zen4-Pro is purpose-built for this regime. While sharing the core Hierarchical Hybrid Attention
architecture with Zen4 (14B), Zen4-Pro's 72B parameter count provides substantially deeper
feed-forward capacity, enabling richer world-model representations that support multi-step
planning. The model's post-training pipeline introduces a dedicated \textbf{Agentic Alignment}
phase that teaches plan-execute-verify patterns from real engineering task traces, not synthetic
dialogues.

The result is a model that scores 48.7\% on SWE-bench Verified --- the industry-standard
evaluation for autonomous software engineering --- while matching or exceeding models of greater
scale on standard reasoning and knowledge benchmarks.

\subsection{Key Contributions}

\begin{itemize}
    \item \textbf{Agentic Alignment Training}: A new post-training stage using real engineering
    task traces with ground-truth verifiable outcomes as the reward signal.
    \item \textbf{Hierarchical Planning}: An internal representation of task state that persists
    across tool calls, enabling reliable multi-step execution without prompt re-injection.
    \item \textbf{Self-Verification}: The model learns to check its own intermediate outputs
    against stated task requirements before proceeding, reducing error propagation in long chains.
    \item \textbf{Frontier Benchmark Results}: New state-of-the-art at 72B dense for SWE-bench,
    GPQA, and MATH.
\end{itemize}

\section{Architecture}

\subsection{Model Specifications}

\begin{table}[H]
\centering
\begin{tabular}{ll}
\toprule
\textbf{Component} & \textbf{Specification} \\
\midrule
Total Parameters         & 72.7B \\
Active Parameters        & 72.7B (dense) \\
Architecture             & Zen MoDE Hybrid Transformer \\
Attention Mechanism      & Hierarchical Hybrid Attention (HHA) \\
Local Window Size        & 8,192 tokens \\
Global Anchor Tokens     & 1,024 per layer \\
Context Length           & 131,072 tokens (128K) \\
Hidden Dimension         & 8,192 \\
Intermediate Dimension   & 29,568 \\
Attention Heads          & 64 \\
Key-Value Heads          & 8 (GQA) \\
Layers                   & 80 \\
Vocabulary Size          & 151,936 \\
Positional Encoding      & RoPE ($\theta = 10{,}000{,}000$) \\
Normalization            & RMSNorm \\
Activation               & SwiGLU \\
\bottomrule
\end{tabular}
\caption{Zen4-Pro Architecture Specifications}
\end{table}

\subsection{Hierarchical Hybrid Attention at 72B}

Zen4-Pro scales the HHA mechanism from Zen4 with a larger local window ($w = 8192$ vs 4096)
and more global anchor tokens ($k = 1024$ vs 512). The larger window improves performance on
tasks requiring dense local reasoning, such as code editing within a single large file.
The additional anchors improve global coherence for tasks that span many documents.

At 72B parameters, the feed-forward networks can represent substantially richer world-model
priors. Ablation studies show that the improvement from Zen4 to Zen4-Pro on complex planning
tasks ($+18\%$ on AgentBench) is primarily attributable to feed-forward capacity rather than
attention mechanism differences, confirming that scale continues to benefit planning tasks even
after architectural improvements have maximized attention efficiency.

\subsection{Agentic State Representation}

Zen4-Pro introduces an explicit task-state representation within the context window. During
agentic post-training, the model learns to maintain a structured scratchpad following the
format:

\begin{lstlisting}[caption=Agentic State Token Format]
<task_state>
  <goal>Refactor authentication module to use JWT</goal>
  <completed>
    - Read auth/session.py (lines 1-340)
    - Identified 3 session-creation call sites
    - Updated auth/session.py with JWT logic
  </completed>
  <pending>
    - Update call sites in api/routes.py, api/middleware.py
    - Run test suite to verify no regressions
  </pending>
  <open_questions>
    - Token expiry: use 24h or match existing session TTL?
  </open_questions>
</task_state>
\end{lstlisting}

This structured scratchpad is generated and updated by the model itself; it is not injected
externally. The model learns to parse its own prior state tokens as factual context, enabling
reliable task tracking over 20+ sequential tool calls without context-window management by
the orchestration layer.

\subsection{Zen MoDE Distillation at Professional Scale}

Zen4-Pro is distilled from an ensemble of teacher models using the fourth-generation Zen MoDE
framework. At 72B, the student has sufficient capacity to absorb richer intermediate
representations, enabling distillation from multiple specialized teachers simultaneously:

\begin{itemize}
    \item A mathematics specialist teacher contributes to layers 20--40.
    \item A code reasoning teacher contributes to layers 40--60.
    \item A general instruction-following teacher provides global signal across all layers.
\end{itemize}

Multi-teacher distillation with layer-targeted routing improves MATH by 3.1\% and HumanEval
by 2.8\% over single-teacher distillation at matched compute.

\section{Training}

\subsection{Pretraining}

\begin{table}[H]
\centering
\begin{tabular}{lcc}
\toprule
\textbf{Data Category} & \textbf{Proportion} & \textbf{Tokens} \\
\midrule
Web (quality-filtered)         & 38\%  & 7.6T \\
Code (multi-language)          & 32\%  & 6.4T \\
Scientific and technical       & 14\%  & 2.8T \\
Tool-call trajectories         & 8\%   & 1.6T \\
Mathematical content           & 5\%   & 1.0T \\
Books and long-form prose      & 3\%   & 0.6T \\
\midrule
\textbf{Total}                 & 100\% & 20.0T \\
\bottomrule
\end{tabular}
\caption{Zen4-Pro Pretraining Data Composition (20T tokens)}
\end{table}

Zen4-Pro pretrained on 20 trillion tokens. The increased code proportion (32\% vs 28\% in Zen4)
reflects the professional-tier emphasis on software engineering tasks. Scientific and technical
text proportion is also increased (14\% vs 10\%) for the deep-reasoning requirements of
professional domains.

\subsection{Post-Training Pipeline}

\subsubsection{Supervised Fine-Tuning}

The SFT stage for Zen4-Pro uses 4.5M instruction-response pairs, approximately twice the
Zen4 volume. The additional data is concentrated in:
\begin{itemize}
    \item Multi-step software engineering tasks with execution traces.
    \item Multi-document analysis and synthesis tasks.
    \item Long-horizon planning and decomposition tasks.
\end{itemize}

\subsubsection{Agentic Alignment}

The agentic alignment stage is unique to Zen4-Pro and Zen4-Max. It trains the model on
trajectories collected from real engineering tasks in sandboxed environments:

\begin{enumerate}
    \item \textbf{Trajectory Collection}: Deploy an early checkpoint on 50,000 real GitHub
    issues in isolated container environments. Collect all tool calls, intermediate states,
    and final outcomes (pass/fail against the issue's test suite).
    \item \textbf{Outcome Filtering}: Retain trajectories where the final outcome is
    verifiably correct (test suite passes). Discard partial successes.
    \item \textbf{GRPO Training}: Use Group Relative Policy Optimization with the binary
    test-pass signal as reward. Generate $G = 8$ trajectories per task and train on the
    relative reward within each group.
    \item \textbf{Iteration}: Repeat trajectory collection and training for 3 rounds
    using the updated model.
\end{enumerate}

This iterative process yields a 14.2-point improvement on SWE-bench Verified over the
SFT-only baseline.

\subsubsection{Preference Optimization}

DPO training uses 1.2M comparison pairs. For professional-tier use cases, comparison pairs
are weighted to emphasize:
\begin{itemize}
    \item Correctness of multi-step plans (3$\times$ weight).
    \item Appropriate tool selection and argument accuracy (2$\times$ weight).
    \item Response clarity and actionability (1$\times$ weight).
\end{itemize}

\subsection{Training Infrastructure}

\begin{table}[H]
\centering
\begin{tabular}{ll}
\toprule
\textbf{Parameter} & \textbf{Value} \\
\midrule
Hardware           & 4,096 $\times$ H100 80GB SXM \\
Parallelism        & TP=8, PP=8, DP=64 \\
Batch Size         & 8M tokens (global) \\
Peak Learning Rate & $1.5 \times 10^{-4}$ \\
LR Schedule        & Cosine with warmup (3000 steps) \\
Optimizer          & AdamW ($\beta_1=0.9, \beta_2=0.95$) \\
Training Duration  & 32 days (pretraining) \\
\bottomrule
\end{tabular}
\caption{Zen4-Pro Training Infrastructure}
\end{table}

\section{Evaluation}

\subsection{Standard Benchmarks}

\begin{table}[H]
\centering
\begin{tabular}{lccc}
\toprule
\textbf{Benchmark} & \textbf{Zen3-Pro} & \textbf{Zen4} & \textbf{Zen4-Pro} \\
\midrule
MMLU (5-shot)            & 85.7\% & 85.4\% & 91.8\% \\
MATH (4-shot)            & 82.1\% & 81.2\% & 87.6\% \\
HumanEval (0-shot)       & 86.4\% & 88.3\% & 93.2\% \\
GPQA Diamond (0-shot)    & 65.8\% & 67.4\% & 74.3\% \\
IFEval (prompt-strict)   & 80.2\% & 83.1\% & 88.4\% \\
GSM8K (8-shot)           & 92.3\% & 93.7\% & 96.8\% \\
BBH (3-shot)             & 81.4\% & 83.9\% & 89.7\% \\
\bottomrule
\end{tabular}
\caption{Benchmark Comparison Across Generations and Tiers}
\end{table}

\subsection{Agentic Benchmarks}

\begin{table}[H]
\centering
\begin{tabular}{lccc}
\toprule
\textbf{Benchmark} & \textbf{Zen3-Pro} & \textbf{Zen4} & \textbf{Zen4-Pro} \\
\midrule
SWE-bench Verified       & 32.1\% & 34.8\% & 48.7\% \\
AgentBench (overall)     & 58.3\% & 63.1\% & 76.4\% \\
WebArena (overall)       & 31.4\% & 35.2\% & 44.8\% \\
ToolBench (success rate) & 67.2\% & 79.4\% & 88.1\% \\
GAIA (Level 1)           & 71.3\% & 74.8\% & 83.6\% \\
GAIA (Level 2)           & 48.9\% & 52.3\% & 63.7\% \\
\bottomrule
\end{tabular}
\caption{Agentic Capability Benchmarks}
\end{table}

\subsection{Multi-Step Reasoning}

\begin{table}[H]
\centering
\begin{tabular}{lcc}
\toprule
\textbf{Task} & \textbf{Steps} & \textbf{Zen4-Pro Score} \\
\midrule
Plan + execute (3-step)   & 3  & 89.4\% \\
Plan + execute (5-step)   & 5  & 81.2\% \\
Plan + execute (10-step)  & 10 & 68.7\% \\
Self-verify + correct     & 2  & 77.3\% \\
Error recovery after tool failure & --- & 71.8\% \\
\bottomrule
\end{tabular}
\caption{Multi-Step Task Completion Rates (AgentBench decomposition)}
\end{table}

\subsection{Professional Domain Evaluation}

\begin{table}[H]
\centering
\begin{tabular}{lcc}
\toprule
\textbf{Domain} & \textbf{Task} & \textbf{Score} \\
\midrule
Software Engineering & SWE-bench Verified       & 48.7\% \\
Data Analysis        & BIRD SQL (accuracy)      & 73.4\% \\
Scientific Research  & GPQA Diamond             & 74.3\% \\
Legal               & LegalBench (avg)          & 68.1\% \\
Finance             & FinanceBench (accuracy)   & 84.2\% \\
Medicine            & MedQA USMLE (4-option)    & 86.3\% \\
\bottomrule
\end{tabular}
\caption{Professional Domain Evaluation}
\end{table}

\section{Related Work}

Scaling dense transformers beyond 70B parameters has followed a predictable capability-compute
curve \cite{hoffmann2022chinchilla}. Zen4-Pro's 72B represents the practical ceiling for
single-node inference on current data-center GPUs while pushing capability to the frontier.

Agentic AI systems built on language models have grown rapidly since ReAct \cite{yao2023react}
demonstrated that interleaving reasoning traces with action invocations improves task success
on multi-step benchmarks. Subsequent systems \cite{wang2024executable} extended this to
hierarchical planning. Zen4-Pro internalizes these patterns rather than relying on external
orchestration frameworks, reducing latency and improving coherence across long trajectories.

SWE-bench \cite{jimenez2024swebench} has emerged as the canonical evaluation for autonomous
software engineering capability. Current frontier scores range from 30--50\% depending on
scaffolding; Zen4-Pro achieves 48.7\% with a minimal scaffold, demonstrating that model
quality rather than scaffolding complexity drives performance.

\section{Deployment}

\subsection{Inference Configuration}

\begin{table}[H]
\centering
\begin{tabular}{llll}
\toprule
\textbf{Format} & \textbf{VRAM} & \textbf{Context} & \textbf{Hardware} \\
\midrule
BF16             & 145 GB & 128K  & H100 $\times 2$ \\
FP8              & 72 GB  & 128K  & H100 $\times 1$ \\
Q4\_K\_M         & 42 GB  & 32K   & RTX 6000 Ada $\times 1$ \\
Q4\_K\_M         & 48 GB  & 64K   & A6000 $\times 1$ \\
\bottomrule
\end{tabular}
\caption{Zen4-Pro Deployment Hardware Requirements}
\end{table}

\subsection{Recommended Agentic Configuration}

\begin{lstlisting}[language=Python, caption=Zen4-Pro Agentic Task Configuration]
from hanzo import HanzoAI

client = HanzoAI()

# Agentic tasks: enable longer generation, structured output
response = client.chat.completions.create(
    model="zen4-pro-72b-instruct",
    messages=[
        {"role": "system", "content": "You are an autonomous engineering agent."},
        {"role": "user", "content": task_description}
    ],
    tools=toolset,
    tool_choice="auto",
    max_tokens=8192,          # allow full reasoning chains
    temperature=0.2,          # low temperature for agentic reliability
    extra_body={
        "agentic_mode": True,  # enables internal state tracking
        "max_tool_rounds": 20  # hard cap on tool-call loops
    }
)
\end{lstlisting}

\section{Safety and Alignment}

\begin{table}[H]
\centering
\begin{tabular}{lcc}
\toprule
\textbf{Safety Metric} & \textbf{Zen3-Pro} & \textbf{Zen4-Pro} \\
\midrule
Harmful content refusal        & 95.3\% & 98.1\% \\
Over-refusal rate (benign)     & 6.2\%  & 3.4\%  \\
Jailbreak resistance (AdvBench)& 93.7\% & 97.8\% \\
TruthfulQA                     & 74.8\% & 82.3\% \\
Agentic safety (destructive action rate) & 1.2\% & 0.3\% \\
\bottomrule
\end{tabular}
\caption{Safety Metrics: Zen3-Pro vs. Zen4-Pro}
\end{table}

Agentic deployment introduces safety risks beyond standard language model outputs: a model
executing multi-step tool plans can take irreversible actions (file deletion, API calls with
side effects). Zen4-Pro's post-training includes an explicit \textbf{action consequence
evaluation} objective that teaches the model to identify and flag potentially irreversible
operations before executing them.

\subsection{Limitations}

\begin{itemize}
    \item Multi-step task completion rates decay with plan length: 10-step tasks complete
    at 68.7\% vs 89.4\% for 3-step tasks, indicating compounding errors in long chains.
    \item Professional domain evaluations (legal, medical) reflect the state of training
    data and should not substitute for domain expert review.
    \item 128K context is sufficient for most engineering tasks but falls short of the
    1M context available in Zen4 (14B). For long-document tasks, Zen4 may be preferable
    despite its lower benchmark scores.
\end{itemize}

\section{Conclusion}

Zen4-Pro delivers frontier professional capability at 72B dense parameters through three
interlocking advances: the Hierarchical Hybrid Attention architecture shared across the
Zen4 family, multi-teacher Zen MoDE distillation that concentrates specialist knowledge
into student layers, and the novel Agentic Alignment post-training stage that grounds
multi-step planning in verifiable engineering outcomes. With 48.7\% on SWE-bench Verified
and strong results across mathematics, science, and knowledge benchmarks, Zen4-Pro is the
recommended choice for enterprise deployments requiring reliable autonomous task execution.

\section*{Acknowledgments}

We thank the annotators, engineers, and researchers at Hanzo AI and Zoo Labs Foundation
who contributed to the trajectory collection, post-training pipeline, and evaluation
infrastructure underlying this work.

\bibliographystyle{plain}
\bibliography{references}

\appendix

\section{Model Card}

\begin{table}[H]
\centering
\begin{tabular}{ll}
\toprule
\textbf{Field} & \textbf{Value} \\
\midrule
Model Name      & Zen4-Pro \\
Version         & v2026.01 \\
Release Date    & January 2026 \\
Parameters      & 72.7B \\
Context Length  & 131,072 tokens \\
License         & Apache 2.0 \\
Repository      & \href{https://huggingface.co/zenlm/zen4-pro-72b-instruct}{huggingface.co/zenlm/zen4-pro-72b-instruct} \\
Documentation   & \href{https://papers.zenlm.org/zen4-pro}{papers.zenlm.org/zen4-pro} \\
Contact         & research@hanzo.ai \\
\bottomrule
\end{tabular}
\caption{Zen4-Pro Model Card}
\end{table}

\end{document}
